\begin{thebibliography}{99}
\bibitem[Andrews et al.(2026)]{acl2026}
Andrews, R., Connes, A., \& La Gatta, A. (2026). The uncertainty
operator, the Weil ground state and the zeros of zeta: A numerical
study and conditional positivity of the Weil quadratic form.
\url{https://arxiv.org/abs/2611.13129}.
\bibitem[Berthelot et al.(1971)]{sga6}
Berthelot, P., Grothendieck, A., \& Illusie, L. (Eds.). (1971).
\emph{Th\'eorie des intersections et th\'eor\`eme de Riemann--Roch
(SGA 6)}. Lecture Notes in Mathematics, 225. Springer.
Expos\'e XIV, \S6, pp.~682--685.
\bibitem[Bhargava(2026a)]{bhargava2026roots}
Bhargava, M. (2026a, October 12). Every nonsquare integer not equal
to $-1$ is a primitive root modulo a positive lower density of primes.
\url{https://hexagonmath.org/2610.00228v1}.
\bibitem[Bhargava(2026b)]{bhargava2026points}
Bhargava, M. (2026b, October 16). Rational points on elliptic curves
that are primitive modulo a positive proportion of primes.
\url{https://hexagonmath.org/2610.00294v1}.
\bibitem[Bhargava(2026c)]{bhargava2026cyclic}
Bhargava, M. (2026c, October 24). Cyclic reductions of elliptic curves
at a positive proportion of primes. \url{https://hexagonmath.org/2610.00699v1}.
\bibitem[Bondarenko et al.(2023)]{bondarenko2023}
Bondarenko, A., Radchenko, D., \& Seip, K. (2023).
Fourier interpolation with zeros of zeta and $L$-functions.
\emph{Constructive Approximation, 57}, 405--461.
\url{https://doi.org/10.1007/s00365-022-09599-w}.
\bibitem[Booker et al.(2019)]{booker2019}
Booker, A. R., Milinovich, M. B., \& Ng, N. (2019).
Subconvexity for modular form $L$-functions in the $t$ aspect.
\emph{Advances in Mathematics, 341}, 299--335.
\url{https://doi.org/10.1016/j.aim.2018.10.037}.
\bibitem[Borger(2009)]{borger2009}
Borger, J. (2009). $\Lambda$-rings and the field with one element.
\url{https://arxiv.org/abs/0906.3146}.
\bibitem[Boulsane and Karoui(2017)]{boulsane2017}
Boulsane, M., \& Karoui, A. (2017). The finite Hankel transform
operator: Some explicit and local estimates of the eigenfunctions
and eigenvalues decay rates. \url{https://arxiv.org/abs/1701.04622}.
\bibitem[Burnol(2001)]{burnol2000}
Burnol, J.-F. (2001). An adelic causality problem related to abelian
$L$-functions. \emph{Journal of Number Theory, 87}(2), 253--269.
\url{https://arxiv.org/abs/math/0001013}.
\bibitem[Chandrasekharan and Narasimhan(1964)]{chandrasekharannar1964}
Chandrasekharan, K., \& Narasimhan, R. (1964).
On the mean value of the error term for a class of arithmetical functions.
\emph{Acta Mathematica, 112}, 41--67.
\url{https://doi.org/10.1007/BF02391764}.
\bibitem[Clunie and Kuijlaars(1994)]{cluniekuijlaars1994}
Clunie, J. G., \& Kuijlaars, A. B. J. (1994).
Approximation by polynomials with restricted zeros.
\emph{Journal of Approximation Theory, 79}(1), 109--124.
\url{https://doi.org/10.1006/jath.1994.1116}.
\bibitem[Connes and Consani(2021)]{cc2020}
Connes, A., \& Consani, C. (2021). Weil positivity and trace formula,
the archimedean place. \emph{Selecta Mathematica, 27}, Article 77.
\url{https://arxiv.org/abs/2006.13771}.
\bibitem[Connes and Consani(2023)]{cc2021}
Connes, A., \& Consani, C. (2023). Spectral triples and zeta-cycles.
\emph{L'Enseignement Math\'ematique, 69}(1--2), 93--148.
\url{https://arxiv.org/abs/2106.01715}.
\bibitem[Connes and Moscovici(2021)]{cm2021}
Connes, A., \& Moscovici, H. (2021). Prolate spheroidal operator and
Zeta. \url{https://arxiv.org/abs/2112.05500}.
\bibitem[Connes et al.(2007)]{ccm2007}
Connes, A., Consani, C., \& Marcolli, M. (2007). The Weil proof and
the geometry of the adeles class space. \url{https://arxiv.org/abs/math/0703392}.
\bibitem[Connes et al.(2024)]{ccm2023}
Connes, A., Consani, C., \& Moscovici, H. (2024). Zeta zeros and
prolate wave operators: Semilocal adelic operators.
\emph{Annals of Functional Analysis, 15}(4), Article 87.
\url{https://arxiv.org/abs/2310.18423}.
\bibitem[Connes et al.(2025)]{ccm2025}
Connes, A., Consani, C., \& Moscovici, H. (2025). Zeta spectral triples.
\url{https://arxiv.org/abs/2511.22755}.
\bibitem[Connes(1999)]{connes1999}
Connes, A. (1999). Trace formula in noncommutative geometry and the
zeros of the Riemann zeta function. \emph{Selecta Mathematica, 5},
29--106. \url{https://arxiv.org/abs/math/9811068}.
\bibitem[Deligne(1974)]{deligne1974}
Deligne, P. (1974). La conjecture de Weil: I.
\emph{Publications Math\'ematiques de l'IH\'ES, 43}, 273--307.
\url{https://doi.org/10.1007/BF02684373}.
\bibitem[Deligne and Serre(1974)]{deligneserre1974}
Deligne, P., \& Serre, J.-P. (1974). Formes modulaires de poids 1.
\emph{Annales scientifiques de l'\'Ecole Normale Sup\'erieure,
S\'erie 4, 7}(4), 507--530.
\url{https://doi.org/10.24033/asens.1277}.
\bibitem[Hoffstein and Ramakrishnan(1995)]{hoffsteinram1995}
Hoffstein, J., \& Ramakrishnan, D. (1995). Siegel zeros and cusp forms.
\emph{International Mathematics Research Notices, 1995}(6), 279--308.
\url{https://doi.org/10.1155/S1073792895000225}.
\bibitem[Horawa(n.d.)]{horawaLfunctions}
Horawa, A. (n.d.). \emph{L-functions} [Lecture notes].
\url{https://people.maths.ox.ac.uk/horawa/L-functions.pdf}.
Page locators use the printed pagination.
\bibitem[Hu(2026)]{hu2026}
Hu, B. (2026). A seven-eighths zero-free half-plane for unitary Weil
$L$-functions. \url{https://hexagonmath.org/2610.00216v2}.
\bibitem[Humphries(2025)]{humphries2025}
Humphries, P. (2025). Archimedean newform theory for $\mathrm{GL}_n$.
\emph{Journal of the Institute of Mathematics of Jussieu, 24}(1), 41--116.
\url{https://doi.org/10.1017/S1474748024000227}.
\bibitem[Jiang and Li(2024)]{jiangli2024}
Jiang, D., \& Li, Z. (2024). The Voronoi summation formula for
$\mathrm{GL}_n$ and the Godement--Jacquet kernels.
\url{https://arxiv.org/abs/2306.02554v4}.
\bibitem[Jiang and Luo(2023)]{jiangluo2023}
Jiang, D., \& Luo, Z. (2023). Certain Fourier operators and their
associated Poisson summation formulae on $\mathrm{GL}_1$.
\emph{Pacific Journal of Mathematics, 326}(2), 301--372.
\url{https://doi.org/10.2140/pjm.2023.326.301}.
The statement and page locators used here refer to the preprint,
\url{https://arxiv.org/abs/2108.03566v2}.
\bibitem[Khare and Wintenberger(2009)]{kharewintenberger2009}
Khare, C., \& Wintenberger, J.-P. (2009).
Serre's modularity conjecture (I).
\emph{Inventiones mathematicae, 178}(3), 485--504.
\url{https://doi.org/10.1007/s00222-009-0205-7}.
The statement locators used here refer to the author preprint,
\url{https://www.math.ucla.edu/~shekhar/papers/results.pdf}.
\bibitem[Landau and Widom(1980)]{landauwidom1980}
Landau, H. J., \& Widom, H. (1980). Eigenvalue distribution of time
and frequency limiting. \emph{Journal of Mathematical Analysis and
Applications, 77}(2), 469--481.
\bibitem[Landau(1985)]{landau1985}
Landau, H. J. (1985). An overview of time and frequency limiting.
In J. F. Price (Ed.), \emph{Fourier techniques and applications}
(pp. 201--220). Plenum Press.
\url{https://doi.org/10.1007/978-1-4613-2525-3_12}.
\bibitem[Lean FRO and Mathlib Initiative(2026)]{alpogeverification}
Lean FRO and Mathlib Initiative. (2026, November 21).
Lean verification of the Riemann Hypothesis.
\url{https://leanprover-community.github.io/blog/posts/riemann-hypothesis/}.
\bibitem[Li(2012)]{li2012}
Li, X.-J. (2012). Prolate spheroidal wave functions, Sonine spaces,
and the Riemann zeta function. \emph{Journal of Mathematical Analysis
and Applications, 389}(1), 379--393.
\url{https://doi.org/10.1016/j.jmaa.2011.11.068}.
\bibitem[Liu(2026)]{liu2026}
Liu, B. (2026). Mixed character moments and a zero-free half-plane
for unitary Hecke $L$-functions. \url{https://hexagonmath.org/2611.00741v1}.
\bibitem[Matom\"aki et al.(2023)]{mstt2023}
Matom\"aki, K., Shao, X., Tao, T., \& Ter\"av\"ainen, J. (2023).
Higher uniformity of arithmetic functions in short intervals I.
All intervals. \emph{Forum of Mathematics, Pi, 11}, Article e29.
\url{https://arxiv.org/abs/2204.03754}.
\bibitem[Meyer(2005)]{meyer2005}
Meyer, R. (2005). On a representation of the idele class group related
to primes and zeros of $L$-functions. \emph{Duke Mathematical Journal,
127}(3), 519--595. \url{https://doi.org/10.1215/S0012-7094-04-12734-4}.
\bibitem[OpenAI(2026)]{openai2026primitive}
OpenAI. (2026, October 4). Primitive roots for every admissible
integer base. \url{https://github.com/openai/math/tree/main/preprints}.
\bibitem[Petrow and Young(2023)]{petrowyoung2023}
Petrow, I., \& Young, M. P. (2023). The fourth moment of Dirichlet
$L$-functions along a coset and the Weyl bound.
\emph{Duke Mathematical Journal, 172}(10), 1879--1960.
\url{https://doi.org/10.1215/00127094-2022-0069}.
\bibitem[Sutherland(2026)]{mitverification}
Sutherland, A. V. (2026, November 24). Verification files for the
Lean proof of the Riemann Hypothesis. Massachusetts Institute of
Technology. \url{https://math.mit.edu/~drew/grh/}.
\bibitem[Suzuki(2026)]{suzuki2026}
Suzuki, M. (2026). Spectral convergence associated with Weil's
quadratic form. \url{https://arxiv.org/abs/2610.32704}.
\bibitem[Thorner and Zaman(2019)]{thornerzaman2019}
Thorner, J., \& Zaman, A. (2019).
A unified and improved Chebotarev density theorem.
\emph{Algebra \& Number Theory, 13}(5), 1039--1068.
\url{https://doi.org/10.2140/ant.2019.13.1039}.
\bibitem[Zeta and L-functions Workbench(2026a)]{simplicitynotes}
Zeta and L-functions Workbench. (2026a).
\emph{The simplicity question for the Riemann zeros after Alp\"oge's theorem}
[Research manuscript, November 23]. Records at \url{\workbenchurl}.
\bibitem[Zeta and L-functions Workbench(2026b)]{commutatornotes}
Zeta and L-functions Workbench. (2026b).
\emph{The exact rank-one commutator does not supply a shrinking-gap rate}
[Research note, November 23]. Records at \url{\workbenchurl}.
\bibitem[Zeta and L-functions Workbench(2026c)]{cubiccontentnotes}
Zeta and L-functions Workbench. (2026c).
[Research manuscript on the joint-content tail for the cubic Kummer
Hecke function, November 24]. Theorems ``Joint-content tail'' and
``The exact remaining square-root criterion''.
Records at \url{\workbenchurl}.
\end{thebibliography}
